% !TeX root = ../papers/nc-p3.tex
\section{全幾何単元対での七次四中心商と十三個の残余軌道}
\label{sec:rank-fifteen-bivariate-seventh-cut}
\begin{theorem}
標数3の任意の幾何単元対 $X\ne1$, $Y\ne0$ において、四始点すべてで
四次から七次の次元が順に $31,40,44,40$ となる共通の四次中心四本と表示が存在する。
下表の十三個のFrobenius軌道では、四中心商を十三次まで認証できる。
従ってこれらの剰余を持つ通常な標数零の階数15窓は標準Koszul左右AS正則で、
各始点は四つの中心四次式上で自由階数256である。
\end{theorem}
\begin{proof}
分岐水準と $X\in\mathbb F_9$ は既知である。
残りでは第\ref{sec:rank-fifteen-bivariate-sixth-cut}節の四列を使い、
語を右に延ばして七次まで進む。各始点で七次の文字別次元は $(10,10,10,10)$ となる。
四始点の主元分子の除外因子の積を同節と同じ最小公倍式として引き戻すと、
$G(v,c)$ の次数は $(328,280)$ である。$X=c^2$, $B=(v^3-v)^2=Y/X^3$ 上の
十二基底係数のモニック最大公約元は $D=X^3B-X+1$ となる。

まず $D\ne0$ とする。十二係数を $D$ で割り、開集合で単元となる
$X,X-1,B,X^3B-(1-X)^2$ の因子も除く。非零係数を $c$ の指数、次いで $v$ の指数の
昇順に並べ、最初の六係数を順に二本ずつ組にして $B$ を消去する。三終結式から
$X,X-1$ 因子を除いた最大公約元の次数は70である。その各既約因子の根で十二係数を
再度 $B$ の多項式として比較する。$[\mathbb F_3(X):\mathbb F_3]>2$ で
共通根を持つのは下表の十一組であり、全て $B\in\mathbb F_3(X)$ である。

次に $D=0$ とする。始点0,2の消去はそのまま使い、始点1,3で別の主元を選ぶ。
各有理関数行をスカラー倍の原始多項式行に直し、台が最小の行のうち原始係数の次数が
最小の位置を取る。同順位なら列順、行順を使う。行の内容と分母をスカラーで保持し、
記録する主元にはこのスカラーを戻す。従って有理関数行列の通常の消去をそのまま
特殊化できる領域が得られる。二通り目の除外積は次数 $(391,347)$ である。
十二係数に $B=(X-1)/X^3$ を代入し、各係数の分母を払った最大公約元は
\[
 X^{175}(X-1)^{39}(X+1)(X^2+X+2)(X^3+2X+2)(X^3+X^2+X+2)^2 .
\]
次数が二を超える二因子は下表の最後の列に対応する。
\[
\begin{array}{c|l|l}
 X\text{ の最小多項式}&B\ (D\ne0\text{ 側の候補})&B\ (D=0)\\ \hline
 X^3+2X+2 & X^2+X+1 & X^2+2X+1\\
 X^3+X^2+2 & 1,\ X^2+2X+2 & \\
 X^3+X^2+X+2 & 1,\ X & 2X^2+2X+1\\
 X^3+2X^2+2X+2 & 2X+2,\ X^2+2X+2,\ 2X^2+2X+1 & \\
 X^4+X^2+2X+1 & X^3+2X+1 & \\
 X^4+X^3+2X+1 & 2X^3+X^2+X & \\
 X^6+X^2+2X+1 & X^5+X^3+X+1 &
\end{array}
\]
\path{computations/verify_rank_fifteen_bivariate_seventh_cut.py} はこの各組を
有限体上の一つの表示で確かめ、Frobeniusで共役な全点を覆う。
実際に平方根を取り、元の四窓から五次までの経路、四次中心、四中心商の
十三次までの次元を計算する。必要な拡大次数はそれぞれ横座標の次数の二倍で足りる。
よって主元が全表示で失敗しうる残余にも同じ始点別の七次次元が存在する。
表の点では第\ref{sec:rank-fifteen-quadratic-verticals}節の四中心の持上げと
自由加群比較を適用できる。
\end{proof}
